\chapter{Introduction, Health Context and Digital Health as a Systemic Enabler}

Digitalization \& Internet Infrastructure

\section{Introduction: Digital Health as a Systemic Enabler}\label{introduction-digital-health-as-a-systemic-enabler}

Digital transformation is now recognized as a foundational enabler of
health system reform, not merely a technical add-on. In Uzbekistan, the
digital health agenda is explicitly linked to the broader goals of
universal health coverage (UHC), equity, efficiency, and quality
improvement. The government's ``Digital Uzbekistan 2030'' strategy and
the National Health System Strategy 2030 both position digitalization as
a cross-cutting priority, aiming to create robust governance structures,
integrated platforms, and the necessary broadband and ICT infrastructure
to support health system modernization. (The World Bank, 2025)

However, the digital health landscape in Uzbekistan is shaped by the
legacy of Soviet-era information systems, the rapid but uneven expansion
of internet and mobile technologies, and persistent gaps in digital
literacy and organizational readiness. This section provides a
comprehensive analysis of the current state, challenges, and
opportunities for digitalization and internet infrastructure in
Uzbekistan's health sector.

\section[Post-2016 Reforms: Infrastructure, Connectivity, and Policy]{Post-2016 Reforms: Infrastructure, Connectivity, and Policy (The World Bank, 2025)}\label{post-2016-reforms-infrastructure-connectivity-and-policy-the-world-bank-2025}

Since 2016, Uzbekistan has embarked on an ambitious digital
transformation agenda, with health as a priority sector. The ``Digital
Uzbekistan 2030'' strategy sets out to strengthen digital governance,
create integrated information systems, and expand broadband connectivity
across the country. The National Health System Strategy 2030 further
articulates the role of digital health as a driver of system-wide
reform, emphasizing the need for interoperable platforms, data-driven
decision-making, and patient empowerment.

By 2023, up to 81\% of public medical organizations' legal entities had
been provided with internet connectivity. However, the quality and
reliability of internet access vary significantly, especially in rural
areas. The phased installation of local area networks (LANs) in
healthcare facilities is ongoing. Despite these advances, more than 50\%
of public medical organizations still lack LANs beyond administrative
offices, and the estimated need for computers ranges from 30,000 to
65,000 units.

The lack of adequate ICT infrastructure is a critical bottleneck for
digital health. Many facilities are underequipped, with insufficient
computers, outdated hardware, and no dedicated ICT maintenance staff.
The absence of robust data centers further limits the scalability of
national digital health systems.

\section{National and Local Health Information Systems}\label{national-and-local-health-information-systems}

Uzbekistan's health sector is characterized by a proliferation of more
than 35 unconnected nationwide IT e-health systems. These systems
operate in silos, with no integration or common master data usage among
Ministry of Health (MoH) information systems. There is no national
patient registry, no unified registry of healthcare providers, and no
standardized nomenclature or classifiers. As a result, integration of IT
systems across facilities is not feasible, and data exchange is severely
constrained. (WHO Regional Office for Europe, 2023)

At the facility level, only a limited number of organizations use
Patient Administration Systems (PAS), Electronic Medical Records (EMR),
or Hospital Information Systems (HIS). Under MoH initiatives, two
PAS/EMR systems have been piloted in Syrdarya region, and about five
local PAS/HIS vendors provide services to public and private clinics and
hospitals. However, these systems are not interoperable, and most
medical records remain paper-based, scattered across facilities, and
inaccessible for longitudinal patient care. (The World Bank, 2025)

\section{Digital Literacy and Adoption Rates (WHO Regional Office for Europe, 2023)}\label{digital-literacy-and-adoption-rates-who-regional-office-for-europe-2023}

Digital literacy among health professionals is highly variable. While
younger staff and those in urban centers are increasingly proficient
with computers and digital tools, many older providers lack basic ICT
skills. Continuous IT training is not yet institutionalized, and there
is no national strategy for digital capacity building in the health
workforce. This gap is compounded by the lack of standardized digital
workflows and the persistence of paper-based documentation, which
consumes significant provider time and undermines efficiency.

Patient digital literacy is also uneven, with urban populations more
likely to use digital health services than rural residents. The absence
of a national patient portal or unified engagement platform limits
opportunities for patients to access their health records, book
appointments, or participate in their care. Where electronic booking
systems exist, their use is limited and not widely spread. As a result,
patients are not ``in the driver's seat'' of their health journey, and
digital health remains largely provider-centric.

The success of digital health initiatives depends not only on technical
infrastructure but also on organizational readiness and change
management. In Uzbekistan, digitalization planning and management are
currently centralized within the MoH and its IT-focused company,
UZINFOCOM. There is no independent body tasked with overseeing digital
health development, and the process is often non-transparent to external
stakeholders, including hospitals, clinics, professional associations,
and patient groups. This lack of inclusive governance hinders buy-in and
slows the pace of digital transformation. (The World Bank, 2025)

